\section{Síntesis y ampliación}

\subsection{Un hilo conductor sistémico que atraviesa toda la clase}

Esta sesión de repaso puede condensarse en una cadena de dependencias: el mundo original se mapea primero mediante tokenización/encoder a representaciones; la arquitectura determina cómo interactúa la información; los datos de preentrenamiento y el objetivo definen qué patrones estadísticos se forman en los parámetros; el post-entrenamiento y el cómputo en tiempo de inferencia modifican el comportamiento; la recuperación, la memoria, las herramientas y el bucle del agente añaden estado externo; la interpretabilidad, el alineamiento y el protocolo de evidencia deciden si el sistema puede entenderse y restringirse; JEPA y SSM exploran respectivamente nuevos objetos de predicción y rutas de cálculo.

Esta cadena también explica por qué las fallas sistémicas suelen propagarse transversalmente entre capas. El tokenizador al romper los límites de las entidades degrada la consulta de recuperación; si esta devuelve evidencia errónea, la trazabilidad del razonamiento se alarga sobre premisas incorrectas; si el modelo de preferencias favorece un tono confiado, el agente podría reducir la abstención necesaria; si la memoria guarda resúmenes no verificados, cada tarea posterior convertirá un error en un estado persistente. Por el contrario, las reparaciones locales también pueden fracasar: ampliar solo el contexto no corrige los permisos erróneos, aplicar solo DPO no actualiza conocimientos obsoletos y añadir solo un verificador no recupera variables inobservables. La auditoría de ingeniería debe ubicar problemas capa por capa a lo largo de representación, datos, objetivo, estado, acción y evaluación, en lugar de atribuir todos los problemas a que "el modelo no es lo bastante inteligente".

Por ello, un informe de sistema confiable debe proporcionar simultáneamente los límites temporales de los datos de entrenamiento, las versiones del modelo y del tokenizador, los métodos de post-entrenamiento, las herramientas y permisos disponibles, el ciclo de vida de la memoria, los conjuntos de evaluación y las métricas de costo. Si faltan estas informaciones, incluso si se reproducen las puntuaciones individuales de un benchmark, no será posible replicar el comportamiento real del sistema; el nombre del modelo no constituye una condición experimental completa.

De igual manera, cualquier afirmación sobre una "alternativa a Transformer" debe especificar bajo qué longitud de secuencia, lote, hardware, presupuesto de entrenamiento y umbrales de calidad se realizó la comparación. Los rankings arquitectónicos fuera del workload carecen de un significado estable; lo verdaderamente transferible son los métodos de medición y las condiciones de contorno.

\begin{importantbox}{Las diez preguntas más transferibles de toda la clase}
\begin{enumerate}
  \item ¿En qué tokens se corta la entrada y qué información se pierde en la entrada?
  \item ¿Qué patrones estadísticos recompensa la función objetivo y cuáles omite?
  \item ¿Cuáles son el origen de los datos, la mezcla, el orden y la frecuencia de actualización?
  \item ¿Dónde está el cuello de botella computacional: parámetros, atención, memoria, recuperación o herramientas?
  \item ¿El post-entrenamiento modifica pesos, cambia la búsqueda de candidatos o añade estado externo?
  \item ¿Cuáles son las observaciones, acciones, permisos, presupuesto y condiciones de parada del agente?
  \item ¿Qué miden las métricas: pérdida, éxito de tareas, preferencia o impacto real?
  \item ¿De dónde provienen los errores: creencias, observación, planificación o proxy de recompensa?
  \item ¿Cómo guarda, actualiza, resuelve conflictos y olvida el sistema?
  \item ¿Bajo condiciones iguales de datos, cómputo y servicio, qué dimensión del Pareto mejora la nueva arquitectura?
\end{enumerate}
\end{importantbox}

\subsection{Tres límites de evidencia}

Primero, los estudios sobre BabyLM, RAG, CGLS, fMRI y alucinación presentados en clase corresponden a diseños experimentales específicos y no pueden extrapolarse como leyes universales desde una única escala o conjunto de datos. Segundo, las demostraciones de aplicación y los benchmarks no sustituyen el costo, los permisos, el desplazamiento de distribución y la validación externa presentes en un flujo de trabajo real. Tercero, conceptos como world model, aprendizaje continuo y alineamiento deben definir primero qué objeto se actualiza y qué protocolo de evaluación aplica; de lo contrario, la misma palabra cubrirá problemas mutuamente distintos.

\subsection{Preguntas abiertas}

Lo siguiente merece seguir cuestionando: ¿es posible obtener cálculo de secuencias largas lineales sin sacrificar la capacidad de acceso aleatorio? ¿Podemos unificar recuperación, actualización de parámetros y conflicto de memoria en un sistema de aprendizaje continuo verificable? ¿Podemos construir un benchmark del agente con diagnóstico estratificado de las creencias del modelo del mundo, el planificador y la acción? ¿Podemos hacer que la supervisión del proceso y la evidencia mecanicista se validen mutuamente en lugar de recompensar conjuntamente un estilo superficial de razonamiento?

\subsection{Lecturas adicionales}

\begin{itemize}
  \item Curso oficial: \url{https://web.stanford.edu/class/cs25/}
  \item Grabaciones oficiales: \url{https://www.youtube.com/watch?v=bHSDPgZYie0}
  \item Baby Scale: \url{https://arxiv.org/abs/2603.29522}
  \item Bilingual BabyLM: \url{https://arxiv.org/abs/2603.29552}
  \item Preentrenamiento--escalado RAG: \url{https://arxiv.org/abs/2604.00715}
  \item Escalado de capas guiado por currículo: \url{https://arxiv.org/abs/2506.11389}
  \item Representación fMRI cruzada de red: \url{https://arxiv.org/abs/2603.00786}
  \item Definición unificada de alucinación: \url{https://arxiv.org/abs/2512.21577}
\end{itemize}

\end{document}
